% Options for packages loaded elsewhere
\PassOptionsToPackage{unicode}{hyperref}
\PassOptionsToPackage{hyphens}{url}
\documentclass[
  a4paper,
]{article}
\usepackage{xcolor}
\usepackage[margin=25mm]{geometry}
\usepackage{amsmath,amssymb}
\setcounter{secnumdepth}{-\maxdimen} % remove section numbering
\usepackage{iftex}
\ifPDFTeX
  \usepackage[T1]{fontenc}
  \usepackage[utf8]{inputenc}
  \usepackage{textcomp} % provide euro and other symbols
\else % if luatex or xetex
  \usepackage{unicode-math} % this also loads fontspec
  \defaultfontfeatures{Scale=MatchLowercase}
  \defaultfontfeatures[\rmfamily]{Ligatures=TeX,Scale=1}
\fi
\usepackage{lmodern}
\ifPDFTeX\else
  % xetex/luatex font selection
\fi
% Use upquote if available, for straight quotes in verbatim environments
\IfFileExists{upquote.sty}{\usepackage{upquote}}{}
\IfFileExists{microtype.sty}{% use microtype if available
  \usepackage[]{microtype}
  \UseMicrotypeSet[protrusion]{basicmath} % disable protrusion for tt fonts
}{}
\makeatletter
\@ifundefined{KOMAClassName}{% if non-KOMA class
  \IfFileExists{parskip.sty}{%
    \usepackage{parskip}
  }{% else
    \setlength{\parindent}{0pt}
    \setlength{\parskip}{6pt plus 2pt minus 1pt}}
}{% if KOMA class
  \KOMAoptions{parskip=half}}
\makeatother
\usepackage{longtable,booktabs,array}
\newcounter{none} % for unnumbered tables
\usepackage{calc} % for calculating minipage widths
% Correct order of tables after \paragraph or \subparagraph
\usepackage{etoolbox}
\makeatletter
\patchcmd\longtable{\par}{\if@noskipsec\mbox{}\fi\par}{}{}
\makeatother
% Allow footnotes in longtable head/foot
\IfFileExists{footnotehyper.sty}{\usepackage{footnotehyper}}{\usepackage{footnote}}
\makesavenoteenv{longtable}
\usepackage{graphicx}
\makeatletter
\newsavebox\pandoc@box
\newcommand*\pandocbounded[1]{% scales image to fit in text height/width
  \sbox\pandoc@box{#1}%
  \Gscale@div\@tempa{\textheight}{\dimexpr\ht\pandoc@box+\dp\pandoc@box\relax}%
  \Gscale@div\@tempb{\linewidth}{\wd\pandoc@box}%
  \ifdim\@tempb\p@<\@tempa\p@\let\@tempa\@tempb\fi% select the smaller of both
  \ifdim\@tempa\p@<\p@\scalebox{\@tempa}{\usebox\pandoc@box}%
  \else\usebox{\pandoc@box}%
  \fi%
}
% Set default figure placement to htbp
\def\fps@figure{htbp}
\makeatother
\setlength{\emergencystretch}{3em} % prevent overfull lines
\providecommand{\tightlist}{%
  \setlength{\itemsep}{0pt}\setlength{\parskip}{0pt}}
\usepackage{bookmark}
\IfFileExists{xurl.sty}{\usepackage{xurl}}{} % add URL line breaks if available
\urlstyle{same}
\hypersetup{
  hidelinks,
  pdfcreator={LaTeX via pandoc}}

\author{}
\date{}

\begin{document}

\section{Generation of Fermionic Structure Is Induced, Autocatalytic,
and Pair-Correlated: Assumptions and Consequences of a Universal
Inelastic Map Acting on Waveforms
Alone}\label{generation-of-fermionic-structure-is-induced-autocatalytic-and-pair-correlated-assumptions-and-consequences-of-a-universal-inelastic-map-acting-on-waveforms-alone}

\textbf{Author:} Noriaki Kihara (WF System Co., Ltd.) \textbf{Date:}
2026-08-05 \textbf{Version:} v1 \textbf{Version DOI:}
\href{https://doi.org/10.5281/zenodo.21808092}{10.5281/zenodo.21808092}
\textbf{Concept DOI:}
\href{https://doi.org/10.5281/zenodo.21808091}{10.5281/zenodo.21808091}

\begin{center}\rule{0.5\linewidth}{0.5pt}\end{center}

\subsection{Abstract}\label{abstract}

For a two-wave relational wave system we design, as an explicit
assumption, a single inelastic interaction operation that \textbf{takes
only the waveforms as input and contains no discriminating branches (no
IF statements)}. Within the assumed class (two-wave, pointwise, cubic,
common-phase invariant, no self-scattering, strength reusing the
existing readout \(R=\sin^2\theta\)), conservation of the
zero-square-sum closure (the axiom) uniquely fixes the relative
coefficient, and the vertex is restricted to the single form
\(\delta a = i R\,(|b|^2 a - b^2 \bar a)\). Measuring the consequences
of this assumed map numerically, we find: (1) purely even-harmonic pumps
generate no odd content identically (a parity theorem --- generation is
necessarily an induced process); (2) over a four-decade sweep of the
seed fraction \(f_0\) the initial growth rate obeys the autocatalytic
ignition law \(\mathrm{rate}=C f_0^2\) with \(C=11.45\) constant to
0.4\%; (3) after ignition there is no runaway --- the system relaxes
oscillatorily to a statistical equilibrium \(f^*=0.4690\), close to the
fermionic-mask phase-space fraction 0.494; (4) although no antiparticle
degrees of freedom are input anywhere, products appear only in the
partner bins dictated by the sum rule (the 485 unpredicted empty bins
remain at machine zero, ratio \(10^{-27}\)), as correlated pairs that
settle the hair (winding-number) ledger (coherence 0.83), rising
simultaneously from exact zero. Furthermore, the uniquified vertex
rewrites exactly as \(\delta a = -2R\,\mathrm{Im}(\bar b a)\,b\),
\(\delta b = +2R\,\mathrm{Im}(\bar b a)\,a\) --- the driving scalar is
the single imaginary part of the relative phase of the two waves, and
the vertex is a pointwise two-channel rotation that is \textbf{exactly
solvable in closed form} (closure and power are conserved identically
pointwise; measured drift \(4\times10^{-14}\) per 3000 collisions). This
paper is not a paper that proves the interaction form correct. We report
with the three layers --- assumption, conditional derivation, and
measurement --- strictly separated, and we measure the value of the
assumption not by correctness but by consistency and productivity (the
same laws are reproduced when embedded in the N-body model). All of the
above shows that from a single waveform map with no particle species, no
antiparticles, and no selection rules as input, a discrete parity class,
autocatalytic generation, statistical saturation, and ledger-closing
conjugate pairs emerge as output-side classifications.

\begin{center}\rule{0.5\linewidth}{0.5pt}\end{center}

\subsection{1. The Character of This Paper --- Separating Assumption,
Derivation, and
Measurement}\label{the-character-of-this-paper-separating-assumption-derivation-and-measurement}

This paper neither derives the generative interaction from the axioms
nor proves that it is correct. What this paper does is three things,
stated separately from the outset:

\begin{itemize}
\tightlist
\item
  \textbf{(a) Assumption}: we posit the form of the inelastic
  interaction as an explicit assumption (Section 4).
\item
  \textbf{(b) Conditional derivation}: within the posited class, we
  derive what the axioms force (Sections 5, 6).
\item
  \textbf{(c) Measurement}: we measure the consequences of the assumed
  interaction in numerical experiments (Sections 8--10).
\end{itemize}

What the reader can verify are the derivations (b) and the measurements
(c); (a) is not an object of verification but the starting point of this
paper. The moment this distinction is blurred, the paper loses its
meaning.

\subsection{2. Claims}\label{claims}

\textbf{First --- the most important design principle of this paper ---
we define the interaction as a single universal operation whose only
inputs are the two waveforms.} Following the anonymity principle, the
interior of the map contains no special-case branches (IF statements)
whatsoever: no branch that discriminates particle species, no case
analysis of ``if boson,'' ``if fermion,'' ``if seed,'' and no threshold
switches. The map receives only the waveforms \((a,b)\) and returns only
waveforms. Even the collision strength is not an externally supplied
parameter but is read out from the waveforms themselves (the reflection
rate \(R=\sin^2\theta\), with \(\theta\) the existing waveform readout).
Consequently, every ``boson,'' ``fermion,'' ``seed,'' and ``partner''
appearing in this paper is a \textbf{readout-side classification} (a
spectral mask), not a \textbf{dynamics-side label}. The map does not
know what it is generating. This design is what makes all subsequent
results nontrivial: the parity selection rule, the autocatalytic
ignition, and the pair structure all appear as consequences of a single
indiscriminate expression applied identically to arbitrary inputs.

\textbf{Second, we enumerate all assumptions.} On top of the existing
universal interaction (the elastic part: a two-channel rotation by the
waveform readout \(\theta\), exactly conserving the combined AB power),
we posit the inelastic extension in the following form: (i) two-wave;
(ii) pointwise (products at each point of the register lattice); (iii)
cubic (lowest nonlinear order); (iv) common-phase invariant; (v) no
self-scattering; (vi) no new coupling constant --- the strength reuses
\(R=\sin^2\theta\). \textbf{All six items are design choices, not
derivations} --- although every item is chosen to be consistent with the
first design principle (waveforms as sole input, no internal branches),
and (vi) is its direct consequence. Each item has a motivation
(anonymity, operational algebra, lowest order), but motivation is not
derivation.

\textbf{Third, within the posited class, Axiom 1 (conservation of the
zero-square-sum closure) uniquifies the form.} For the most general form
under the assumptions, \(\delta a = i(g_1|b|^2 a + g_2 b^2\bar a)\),
closure conservation forces \(g_2=-g_1\). Hence within the assumed class
the vertex is restricted to the single form
\(\delta a = igR(|b|^2 a - b^2\bar a)\). \textbf{This is conditional
uniqueness, not uniqueness of the class itself.} It has moreover turned
out that the uniquified vertex is exactly solvable in closed form as a
pointwise two-channel rotation (Section 5.2) --- closure conservation
becomes a pointwise identity and no numerical integration is needed.
Note also that what closure conservation demands of the strengths is
pair symmetry only; the arithmetic-mean rule for mixed collisions is an
\textbf{additional experimental fact} not derived from closure (the
product form also conserves closure --- we do not blur this
distinction).

\textbf{Fourth, the parity theorem.} Under the assumed vertex, no
odd-harmonic content is ever generated from purely even-harmonic pumps.
Hence \textbf{in this model generation is necessarily an induced
process}; spontaneous seedless generation does not occur. We prove this
as a theorem and verify it by machine. A single expression with no IF
statements nonetheless possesses a strict selection rule --- the
selection rule arises not from branching but from the exponent
arithmetic of the expression.

\textbf{Fifth, measurement of the ignition law.} Over a four-decade
sweep of the seed fraction, \(\mathrm{rate} = C f^2\) with
\(C = 11.45\), constant to 0.4\%. The process is autocatalytic, and the
ignition time scales as \(1/f_0^2\). The coefficient has been
calibrated: the initially measured \(C=10.4\) contained an integration
bias of about 10\% from the midpoint method. A RK4 re-run improved the
zero-closure drift from \(2.9\times10^{-3}\) to \(2.8\times10^{-9}\),
and the law and its constancy were invariant across both integrators
(the history is recorded in Section 11).

\textbf{Sixth, measurement of the post-ignition fate.} Total conversion
(runaway) does not occur. The vertex is reversible, and the system
oscillates back and forth to a statistical equilibrium \(f^* = 0.4690\).
This value nearly coincides with the phase-space fraction of the
fermionic mask, 0.494, and reads as equipartition thermalization.

\textbf{Seventh, measurement of the pair structure (three census
criteria).} Although antiparticle-like degrees of freedom are input
neither in the initial conditions nor in the map: (1) generation appears
only in the partner bins dictated by the sum rule, and the 485
unpredicted empty bins remain at machine zero (\(10^{-27}\) of the
partner-band mean); (2) the partner shows coherence 0.83 with the seed
at the position \(q^*=+4\) predicted by the hair sum rule (the ledger is
machine-derived from the output of a single vertex application ---
nothing hand-set); (3) the partner band rises from exact zero
simultaneously with the sidebands. Conclusion: \textbf{inside this
model, antiparticles are outputs, not inputs.}

\textbf{Eighth, we record the full process of refutation and correction
(Section 11).}

\textbf{Finally, what this paper does not claim, and the status of the
assumption (Sections 13, 14).} This paper does not claim that nature's
generative interaction has this form. What raises the value of the
assumption is not a proof of correctness but a demonstration of
consistency and productivity. The assumption is falsifiable.

\subsection{3. Notation and Foundation --- the Elastic Universal
Interaction}\label{notation-and-foundation-the-elastic-universal-interaction}

The foundation is the two-wave relational wave system of Paper 9
{[}1{]}. The state is a pair of complex waveforms \((a,b)\) (a \(\chi\)
lattice of 512 bins \(\times\) an \(\eta\) lattice of 16 points). The
existing \textbf{elastic} universal interaction applies, at each
collision,

\[\theta = \operatorname{atan2}(\sqrt{P_f}, \sqrt{P_b}), \qquad R = \sin^2\theta,\]

\[\begin{pmatrix} a' \\ b' \end{pmatrix} = \begin{pmatrix} \cos\theta & -\sin\theta \\ \sin\theta & \cos\theta \end{pmatrix}\begin{pmatrix} a \\ b \end{pmatrix}\]

where \(P_f, P_b\) are the fermionic/bosonic relational quantities of
the combined AB spectrum (the fermionic mask = the bin set with even
\(|k|\ge 4\)). Hereafter we call \(f = P_f/\lVert(a,b)\rVert^2\), the
ratio of fermionic relational power to total power, the fermionic
fraction. This rotation exactly conserves the per-bin power of the AB
combination --- it is \textbf{elastic} and has no channel that changes
harmonic content. Hence generating fermionic structure from a bosonic
sea requires an added inelastic channel. This is the question from which
this paper starts.

The readout \(\theta\) (and hence \(R\)) is determined by the waveforms
alone. The construction and verification of this readout belong to Paper
9 {[}1{]}; in this paper the original code is referenced read-only as an
unmodified snapshot (SHA-256 recorded; Section 15).

\subsection{4. Assumptions --- Design of the Inelastic
Extension}\label{assumptions-design-of-the-inelastic-extension}

\subsubsection{4.1 Design principle (Claim
One)}\label{design-principle-claim-one}

The inelastic extension must be a \textbf{universal operation} with the
same standing as the elastic part: inputs are the waveforms \((a,b)\)
only, the output is waveforms only, and there is no internal branch
discriminating particle species or situations. The only conditional in
the implementation is the numerical skip ``do nothing if the strength is
zero'' (\texttt{if\ r\ \textgreater{}\ 0}), which is not a decision rule
--- the strength \(r=R\) is itself a readout of the waveforms. This was
confirmed by source audit (test T5).

\subsubsection{4.2 Assumption list}\label{assumption-list}

{\def\LTcaptype{none} % do not increment counter
\begin{longtable}[]{@{}
  >{\raggedright\arraybackslash}p{(\linewidth - 4\tabcolsep) * \real{0.1034}}
  >{\raggedright\arraybackslash}p{(\linewidth - 4\tabcolsep) * \real{0.2069}}
  >{\raggedright\arraybackslash}p{(\linewidth - 4\tabcolsep) * \real{0.6897}}@{}}
\toprule\noalign{}
\begin{minipage}[b]{\linewidth}\raggedright
\#
\end{minipage} & \begin{minipage}[b]{\linewidth}\raggedright
Assumption
\end{minipage} & \begin{minipage}[b]{\linewidth}\raggedright
Motivation (not derivation)
\end{minipage} \\
\midrule\noalign{}
\endhead
\bottomrule\noalign{}
\endlastfoot
(i) & Two-wave (no vertex of three or more waves) & A collision is a
two-wave event \\
(ii) & Pointwise (products at each point of the \(\chi\times\eta\)
lattice) & Four-wave-mixing sum rules of type \(k_1+k_2-k_3=k_4\) hold
automatically as lattice arithmetic; translation invariant \\
(iii) & Cubic (lowest nonlinear order) & Vanishes in the dilute limit,
reducing to the elastic operation (correspondence principle) \\
(iv) & Common-phase invariant & The convention that the overall phase is
meaningless \\
(v) & No self-scattering (no \(|a|^2a\)-type term) & Operational
algebra: ``a single wave has no self-scattering'' \\
(vi) & No new coupling constant; strength \(=R=\sin^2\theta\) &
Anonymity: introduce no new free parameter. The complement of \(R\) (the
transmission rate \(1-R\)) has been identified with the elementary
charge {[}4{]} --- the strength is a readout quantity that already has
physical standing \\
\end{longtable}
}

These six items are design choices. The motivations make each item
\textbf{plausible}, but do not \textbf{force} it.

\subsubsection{4.3 Implemented form}\label{implemented-form}

One collision = elastic part (unchanged) \(\circ\) inelastic part:

\[a' = a + i\,R\,(|b|^2 a - b^2 \bar a), \qquad b' = b + i\,R\,(|a|^2 b - a^2 \bar b)\]

(simultaneous update, pointwise; the relative sign is fixed by the
theorem of the next section. The numerical implementation uses the
closed-form exact solution of Section 5.2 --- the RK4 substeps are the
historical implementation of the calibration process, recorded in
Section 11.)

\subsection{5. Theorem 1: Closure Conservation Uniquifies the
Coefficient}\label{theorem-1-closure-conservation-uniquifies-the-coefficient}

The axiom is conservation of the zero-square-sum closure
\(\mathcal{C} = \sum a^2 + \sum b^2 = 0\). For the most general vertex
under assumptions (i)--(v),

\[\delta a = i(g_1 |b|^2 a + g_2 b^2 \bar a), \qquad \delta b = i(g_1 |a|^2 b + g_2 a^2 \bar b),\]

the first variation is

\[\delta \mathcal{C} = 2\sum (a\,\delta a + b\,\delta b) = 2i\,(g_1+g_2)\sum\left(|b|^2 a^2 + |a|^2 b^2\right).\]

The sum on the right does not vanish in general, so
\(\delta\mathcal{C}=0 \iff g_2 = -g_1\). Hence \textbf{within the
assumed class} the vertex is restricted to the single form of Section
4.3 (the strength absorbs \(g_1\) into \(R\) by redefinition).

We make three points explicit. First, this is \textbf{conditional
uniqueness}: with the class (i)--(v) posited, the axiom selects the
coefficient; the class itself has not been derived. Second, what closure
conservation demands of the strengths is \textbf{pair symmetry only}. As
the mixing rule when the two waves have different strengths, both the
arithmetic mean and the product conserve closure exactly (numerically
verified; T5 of {[}9{]}). That the actual dynamics is consistent with
the arithmetic mean and inconsistent with the product is an
\textbf{additional experimental fact} not derived from closure. Third,
as a historical record, the first version v1 was implemented in the
general form (\(g_2=+g_1\)) and closure drift was observed --- the
current form is the correction to \(g_2=-g_1\) demanded by closure
conservation (Section 11).

\subsubsection{5.1 Corollary 1: the vertex is an exchange flow driven by
the relative
phase}\label{corollary-1-the-vertex-is-an-exchange-flow-driven-by-the-relative-phase}

By an identity transformation (verified to machine precision), the
uniquified vertex can be written

\[\delta a = -2R\,\mathrm{Im}(\bar b a)\,b, \qquad \delta b = +2R\,\mathrm{Im}(\bar b a)\,a.\]

The substance of what looked like a cubic nonlinearity is \textbf{an
exchange flow that reads the single quantity \(s=\mathrm{Im}(\bar b a)\)
--- the imaginary part of the relative phase of the two waves --- and
transports amplitude toward the other wave}. Under the common phase
\((a,b)\mapsto(e^{i\phi}a, e^{i\phi}b)\), \(s\) is invariant, so the
common-phase invariance of assumption (iv) can be read off explicitly
from this form --- the scalar driving the interaction is reduced to one
antisymmetric relative-phase quantity, and the common phase has no way
to enter the expression from the start.

\subsubsection{5.2 Corollary 2: the vertex is a pointwise two-channel
rotation, exactly solvable in closed
form}\label{corollary-2-the-vertex-is-a-pointwise-two-channel-rotation-exactly-solvable-in-closed-form}

The flow of Corollary 1 is a flow of the two-component vector \((a,b)\)
along the tangent direction \((-b, a)\) --- that is, a \textbf{pointwise
two-channel rotation} (angular velocity \(2s\)). Moreover, that \(s\)
itself is exactly invariant under this flow can be shown by direct
computation: for the flow \(\dot a = -2s\,b\), \(\dot b = +2s\,a\)
(\(s\) real),

\[\frac{d}{d\tau}(\bar b a) = \dot{\bar b}\,a + \bar b\,\dot a = 2s\,|a|^2 - 2s\,|b|^2\]

and the right-hand side is real. Hence
\(\dfrac{ds}{d\tau} = \mathrm{Im}\,\dfrac{d}{d\tau}(\bar b a) = 0\).
Since the angular velocity is conserved along the flow, the inelastic
part has a \textbf{closed-form exact solution} at each lattice point:

\[a' = \cos\varphi\,a - \sin\varphi\,b, \qquad b' = \sin\varphi\,a + \cos\varphi\,b, \qquad \varphi(x) = 2R\,\mathrm{Im}(\bar b a)(x).\]

There are three consequences. First, the closure \(a^2+b^2\) and the
power \(|a|^2+|b|^2\) are \textbf{conserved identically pointwise}
(invariants of a real rotation) --- the geometric meaning of the
uniqueness theorem (Section 5 main text) becomes transparent: ``the
closure-conserving direction is the tangent direction of rotation.''
Second, numerical integration becomes unnecessary: re-running all
experiments of this paper in closed form reproduced everything ---
3000-collision closure drift \(4.1\times10^{-14}\) and norm drift
\(1.4\times10^{-14}\) (machine precision), with \(f^*=0.4690\), the
maximum of \(f\) and its position, and \(C\) (within 0.5\%) all
recovered (criteria E1--E3 fixed in advance; all passed). Third, it
becomes visible that the elastic and inelastic parts share the
\textbf{same rotation-generating direction}: one collision is ``a global
rotation by a finite angle read from the spectrum'' \(\circ\) ``a
pointwise rotation by an angle determined by the local relative phase,''
and the two parts can be read in a unified way as rotations of the same
two-channel relational space. Generation (the appearance of new
harmonics) arises from the rotation angle \(\varphi(x)\) differing from
lattice point to lattice point --- the spectrum changes while the
pointwise combined power never changes at all. This unified reading (the
possibility of deriving both from one generator family) is listed as an
unverified problem in Section 13.1.

\subsection{6. Theorem 2: the Parity Theorem --- Generation Is
Necessarily
Induced}\label{theorem-2-the-parity-theorem-generation-is-necessarily-induced}

\textbf{Theorem}: under the assumed vertex, if \(a\) and \(b\) both
contain only even harmonics, the odd-harmonic content remains
identically zero.

\textbf{Sketch of proof}: the vertex is a pointwise cubic product, so in
spectral space it is a convolution composing bins as
\(k_1 + k_2 - k_3\). Even + even - even = even. Hence the even parity
class is closed and no generation term for odd content exists. The
selection rule arises not from branching but from the arithmetic of
exponent sums. \(\blacksquare\)

\textbf{Machine verification} (T3): colliding purely even pumps (no
seed) at strong coupling \(g=0.5\) for 200 collisions, the maximum of
the fermionic relational power stays below \(10^{-25}\) (machine zero).

\textbf{Consequence}: in this model, spontaneous seedless generation
does not occur. Generation is necessarily an \textbf{induced} (seeded)
process --- the coexistence of a tiny odd seed and strong even pumps is
the condition. This is the strict, discrete-map version of the fact that
spontaneous emission does not appear in semiclassical radiation theory
(Section 12). We state explicitly that the spontaneous version (seeded
by zero-point fluctuations) is outside the scope of this map.

\textbf{Positioning}: the even/odd class of harmonic number is the
quotient \(\mathbb{Z}/2\mathbb{Z}\), and the sum rule \(k_1+k_2-k_3\)
conserves it. Parity thus appears not as an externally attached quantum
number but as \textbf{the minimal discrete invariant conserved by the
exponent arithmetic of the interaction}. We emphasize: group theory was
not an input to the design. The accurate order of events is that when
the outputs of the single IF-free expression were classified after the
fact, a conserved class of \(\mathbb{Z}_2\) type was already there.

\subsection{7. Correspondence Principle and the Ledger of Conserved
Quantities}\label{correspondence-principle-and-the-ledger-of-conserved-quantities}

\begin{itemize}
\tightlist
\item
  \textbf{T1 control}: at zero strength, bitwise identical to the
  original pipeline (the inelastic part vanishes completely).
\item
  \textbf{T3' dilute reduction}: the trajectory difference with/without
  the inelastic part vanishes as \(s^2\) in the scale \(s\) (measured
  ratios 3.95, 4.00 per halving --- correspondence principle: reduction
  to the elastic operation in the dilute limit).
\item
  \textbf{Conserved quantities}: by the closed form of Section 5.2, the
  combined AB norm and the zero closure are conserved
  \textbf{identically pointwise}. Measured (3000 collisions, strong
  occupation \(S=8\), closed-form implementation): zero-closure drift
  \(4.1\times10^{-14}\), norm drift \(1.4\times10^{-14}\) (machine
  precision). The drifts observed with the RK2/RK4 implementations of
  the calibration process (\(2.9\times10^{-3}\) / \(2.8\times10^{-9}\))
  are retroactively confirmed to have been pure integrator error
  (Section 11). The combined \(\chi\) momentum is conserved by the sum
  rule via the translation invariance of the pointwise operation.
\end{itemize}

\subsection{8. Experiment 1: the Autocatalytic Ignition Law rate =
C·f²}\label{experiment-1-the-autocatalytic-ignition-law-rate-cfuxb2}

\textbf{Procedure} (preregistered): inject an odd seed (amplitude ratios
\(10^{-3}, 10^{-2}, 10^{-1}\) --- four decades in the initial seed
fraction \(f_0\)) into two even pumps (scale \(S=8\)), measure
\(\mathrm{rate}_0\) from the initial slope of \(\ln P_f\) over 200
collisions (\(j\le 20\)), and test the constancy of
\(C = \mathrm{rate}_0/f_0^2\).

\textbf{Results} (Figure 1):

{\def\LTcaptype{none} % do not increment counter
\begin{longtable}[]{@{}lll@{}}
\toprule\noalign{}
Seed amplitude ratio & \(f_0\) & \(C=\mathrm{rate}_0/f_0^2\) \\
\midrule\noalign{}
\endhead
\bottomrule\noalign{}
\endlastfoot
\(10^{-3}\) & \(4.77\times10^{-7}\) & 11.453 \\
\(10^{-2}\) & \(4.77\times10^{-5}\) & 11.462 \\
\(10^{-1}\) & \(4.74\times10^{-3}\) & 11.410 \\
\end{longtable}
}

Over four decades of \(f_0\), \(C = 11.45\) with spread 0.4\%. The
growth rate is proportional to the \textbf{square} of the seed fraction
--- an autocatalytic law --- and the ignition time scales as
\(1/f_0^2\): a sea with a small seed fraction does not ignite for
astronomical times.

\begin{figure}
\centering
\pandocbounded{\includegraphics[keepaspectratio,alt={Figure 1 ignition law}]{figs/fig1_ignition_law_v1.png}}
\caption{Figure 1 ignition law}
\end{figure}

\textbf{Figure 1}: (a) The growth trajectories
\((\ln f - \ln f_0)/f_0^2\) collapse onto a single common curve over
four decades of \(f_0\) (the extension of the \(f^2\) law to the whole
trajectory). The dashed line is the initial slope \(C=11.45\). The
common gain curve grows superlinearly with collisions (this
superlinearity is common to all seed fractions and is not explained by
the growth of the seed itself --- identifying the mechanism is recorded
as outside the scope of this paper). The largest seed (green) departing
upward from the common curve in the second half is the direct appearance
of autocatalytic acceleration by the current \(f^2\). (b) Constancy of
\(C\). Gray: before calibration (RK2 integration, systematic bias of
about 10\%; Section 11).

\subsection{9. Experiment 2: the Post-Ignition Fate --- No Runaway, but
Statistical
Equilibrium}\label{experiment-2-the-post-ignition-fate-no-runaway-but-statistical-equilibrium}

\textbf{Procedure} (preregistered): \(S=8\), seed amplitude ratio
\(0.1\), 3000 collisions. Record the zero-closure drift, the norm drift,
and the full series of \(f\). Criteria: (R1) closure drift below 1/100
of the pre-calibration measurement \(2.9\times10^{-3}\); (R2) \(f^*\)
(mean of the last 500) within \(0.469\pm0.02\) with the oscillatory-type
verdict retained.

\textbf{Results} (Figure 2): total conversion (runaway to \(f\to1\))
does not occur. \(f\) rises steeply at \(j\approx450\), passes a maximum
\(0.607\) (\(j=1841\)), oscillates, and settles at the statistical
equilibrium \(f^* = 0.4690\). This value nearly coincides with the
phase-space fraction of the fermionic mask, \(0.494\) (the ratio of mask
bins to all bins), and reads as equipartition thermalization under the
reversible vertex (the attribution of the \textasciitilde5\% difference
between the two values is in Section 13.1). The closure drift is
\(4.1\times10^{-14}\) in the closed-form implementation, norm drift
\(1.4\times10^{-14}\) --- machine precision (the calibration RK4 run
gave \(2.8\times10^{-9}\) / \(8.6\times10^{-11}\); both R1 and R2 hold).

\begin{figure}
\centering
\pandocbounded{\includegraphics[keepaspectratio,alt={Figure 2 fate}]{figs/fig2_fate_v1.png}}
\caption{Figure 2 fate}
\end{figure}

\textbf{Figure 2}: the post-ignition fate. Blue is the RK4 calibration
run; gray (RK2, pre-calibration) overlaps the blue almost completely and
is invisible --- the six-orders-of-magnitude difference in integrator
precision does not change the trajectory or the equilibrium value (the
physics is integrator-independent). Red dashed line = statistical
equilibrium \(f^*=0.4690\); green dotted line = mask phase-space
fraction \(0.494\) (candidate equipartition value).

\subsection{10. Experiment 3: the Pair-Structure Census ---
Antiparticles Are Outputs, Not
Inputs}\label{experiment-3-the-pair-structure-census-antiparticles-are-outputs-not-inputs}

\textbf{Design}: so that the pre-thermalization imprint separates bin by
bin, the pumps are single-winding, narrow-band (harmonics \(30,32,34\)
\(\to\) raw bins \(\{29,31,33,35\}\)) \(\times\) two channels; the seed
is single-winding harmonic \(21\) (raw bin \(22\), amplitude ratio
\(0.2\)); the measurement window is limited to 40 collisions.
\textbf{Preregistered predictions} (fixed and printed before execution):

\begin{itemize}
\tightlist
\item
  Partner band (the pair-generation output) = (pump bins \(\oplus\) pump
  bins) - seed bin = \(\{36,38,\dots,48\}\) (7 bins)
\item
  XPM sidebands (phase modulation, not pairs) = seed bin ± pump
  differences = \(\{16,18,20,24,26,28\}\) (6 bins)
\item
  \textbf{All empty bins other than the above plus the pump-sum band and
  the seed neighborhood (485 bins) are predicted not to generate} (a
  machine-zero prediction)
\item
  Hair (\(\eta\) winding) ledger: the partner's hair is
  \textbf{machine-derived from the increment of a single vertex
  application to the initial state} (hand-setting forbidden). The
  derived value agrees with the sum rule of the pair-generation term
  \(b^2\bar a\), \(m^* = 2m_B - m_s = 2(+2)-(+1) = +3\), and the
  predicted winding offset of the pair correlation is
  \(q^* = m_s + m^* = +4\)
\end{itemize}

\textbf{Criteria and results} (all re-measured with RK4; agreement with
the pre-calibration measurements to 4 digits):

\begin{itemize}
\tightlist
\item
  \textbf{P1 sum-rule exclusivity}: the median growth of the 485
  unpredicted bins, \(8.2\times10^{-32}\), is \(1.4\times10^{-27}\) of
  the partner-band mean growth \(5.8\times10^{-5}\) --- machine zero.
  \textbf{Generation appears only in the predicted bins.}
\item
  \textbf{P2 the hair ledger}: the \(q\)-scan coherence (anomalous pair
  correlation of seed \(\times\) partner) is \(0.832\) at \(q^*=+4\) and
  \(7.5\times10^{-8}\) at the uniform average (\(q=0\)) --- not the
  absence of correlation, but the hair selection rule extinguishing the
  \(q=0\) component. \textbf{The partner is born correlated with the
  seed in the form that settles the hair ledger.}
\item
  \textbf{P3 simultaneity of the pair}: the partner band rises from an
  initial \(1.5\times10^{-28}\) (exact zero), and its 40-collision
  growth is \(0.45\) of the sideband growth --- simultaneous output of
  the same vertex.
\end{itemize}

\begin{figure}
\centering
\pandocbounded{\includegraphics[keepaspectratio,alt={Figure 3 census}]{figs/fig3_census_v1.png}}
\caption{Figure 3 census}
\end{figure}

\textbf{Figure 3}: (a) per-bin power growth (40 collisions). Only the
predicted bins grow (red = partner band, blue = XPM sidebands, green =
pump-sum band), and the 485 unpredicted empty bins (gray, the flat part
on the right) stay pinned at machine zero. The gray skirt on the left is
the initial structure of pumps and seed (not empty bins). (b) The
anomalous pair correlation for each hair winding offset \(q\). Only the
prediction \(q^*=+4\) stands at \(0.83\); the others are extinguished
below \(10^{-8}\) by the selection rule.

\textbf{Conclusion}: antiparticle-like degrees of freedom are input
neither in the initial conditions nor in the map. Nevertheless, the
products appeared at the positions dictated by the sum rule, as
correlated pairs carrying the hair ledger, simultaneously from exact
zero. \textbf{Inside this model, antiparticles are outputs, not inputs.}
More precisely, the non-circular definition this model provides is:
\textbf{a partner (antiparticle candidate) is the conjugate relational
state that the generating vertex co-produces in order to close all
conservation ledgers (the bin sum rule \(k^*=2k_p-k_s\) and the hair sum
rule \(m^*=2m_B-m_s\))}. It is not that a wave viewed alone is a
particle or an antiparticle; the distinction is determined as a role
within the generation event --- the appearance, at the level of
generation phenomena, of this series' design principle that relations
precede attributes.

\subsection{11. Record of Refutations and
Corrections}\label{record-of-refutations-and-corrections}

Following the series convention (record the full process of hypothesis
\(\to\) refutation \(\to\) correction), we list every place where this
study actually went wrong.

{\def\LTcaptype{none} % do not increment counter
\begin{longtable}[]{@{}
  >{\raggedright\arraybackslash}p{(\linewidth - 6\tabcolsep) * \real{0.1429}}
  >{\raggedright\arraybackslash}p{(\linewidth - 6\tabcolsep) * \real{0.2857}}
  >{\raggedright\arraybackslash}p{(\linewidth - 6\tabcolsep) * \real{0.2857}}
  >{\raggedright\arraybackslash}p{(\linewidth - 6\tabcolsep) * \real{0.2857}}@{}}
\toprule\noalign{}
\begin{minipage}[b]{\linewidth}\raggedright
\#
\end{minipage} & \begin{minipage}[b]{\linewidth}\raggedright
Error
\end{minipage} & \begin{minipage}[b]{\linewidth}\raggedright
Discovery
\end{minipage} & \begin{minipage}[b]{\linewidth}\raggedright
Correction
\end{minipage} \\
\midrule\noalign{}
\endhead
\bottomrule\noalign{}
\endlastfoot
1 & v1 vertex implemented in the general form \(g_2=+g_1\) & Measured
zero-closure drift & Derived \(g_2=-g_1\) from the first variation of
closure (Section 5). The error turned into the discovery of the
uniqueness theorem \\
2 & census v1: measured the pair imprint after thermalization (after 800
collisions) & Three criteria FAIL & Moved the measurement window to 40
collisions before thermalization \\
3 & census v2: broadband pumps dispersed the partner over 444 bins;
correlation measured with the \(\eta\)-uniform average & Correlation
exactly zero & Changed to narrow-band single-winding design +
hair-resolved correlation. The ``zero'' was not decorrelation but the
hair selection rule \\
4 & Hair ledger preregistered by hand calculation (assuming \(m_B=-1\))
& P2 FAIL (measured \(q^*=+4\neq\) prediction) & The measured carrier
hair of pump B is \(+2\). Replaced the ledger with machine derivation
from a single vertex application --- what was wrong was the hand-set
ledger on the judging side, not the sum rule of the map \\
5 & Integration bias of RK2 (midpoint) substeps & Closure drift
\(2.9\times10^{-3}\); RK4 control in the many-body implementation & RK4
re-run: drift \(2.8\times10^{-9}\) (six orders better). \(f^*\), fate,
census invariant to 4 digits; only \(C\) calibrated \(10.4\to11.45\) (a
\textasciitilde10\% systematic error). Subsequently the discovery of the
closed-form solution (Section 5.2) made the integrator itself
unnecessary (drift \(10^{-14}\)), finally confirming this family of
drifts as pure integrator error \\
6 & Temporarily claimed ``closure conservation selects the arithmetic
mean of mixed strengths'' & Numerical verification: the product form
also conserves closure exactly & Closure demands pair symmetry only. The
selection of the mean is demoted to an additional experimental fact of
the dynamics (Section 5) \\
\end{longtable}
}

\subsection{12. Prior Work --- Correspondences Found After Independent
Study}\label{prior-work-correspondences-found-after-independent-study}

The vertex of this paper was constructed from an independent axiomatic
study (closure conservation) and was not derived by relying on the
following prior fields. Noticing that the measured results (partner sum
rule, induced-only generation, statistical equilibrium) are isomorphic
to known phenomena, we re-surveyed the literature and found the
following correspondences. The correspondences are supporting evidence
that the assumed class is not physically vacuous; they are not proofs of
correctness.

\begin{enumerate}
\def\labelenumi{\arabic{enumi}.}
\tightlist
\item
  \textbf{The algebraic form of the vertex}: the coupled nonlinear
  Schrödinger equations of birefringent optical fibers {[}5{]} carry
  both the XPM term \(|A_y|^2A_x\) and the coherent-coupling term
  \(A_y^2 A_x^*\) --- the same algebraic class as our
  \(g_1|b|^2a + g_2 b^2\bar a\). In optics, however, the coefficient
  ratio is fixed by the \(\chi^{(3)}\) tensor symmetry of the medium,
  whereas here closure conservation (the axiom) forces \(g_2=-g_1\).
  Same class, different selection principle.
\item
  \textbf{The partner sum rule}: in seeded four-wave mixing {[}6{]} the
  idler is automatically generated at
  \(\omega_i = 2\omega_p - \omega_s\) with phase-conjugate correlation
  --- strictly isomorphic to our \(k^* = 2k_p - k_s\). In optics, phase
  matching is imposed by the dispersion relation, whereas our sum rule
  comes from the exponent arithmetic of the pointwise cubic map. We find
  no counterpart of the hair-inclusive ledger and the machine-zero
  exclusivity (\(10^{-27}\)).
\item
  \textbf{Induced processes only}: in semiclassical radiation theory
  only stimulated emission and absorption appear; spontaneous emission
  does not {[}7{]}. The seed of spontaneous four-wave mixing is quantum
  vacuum noise. Our parity theorem is the discrete-map version of this
  fact --- not an approximate statement but a strict selection rule with
  machine verification.
\item
  \textbf{Fermion generation from a bosonic background}: in cosmological
  preheating, the oscillating inflaton parametrically excites fermions,
  saturating under Pauli blocking {[}8{]}. The phenomenology is parallel
  but the ontology differs: preheating posits a Dirac field as input. We
  input no fermionic degrees of freedom --- ``fermion'' is a readout
  mask. The saturation mechanisms also differ (Pauli blocking
  vs.~equipartition equilibrium of a reversible vertex).
\item
  \textbf{Oscillation and thermalization}: FPU-Tsingou recurrence
  {[}10{]} and Rayleigh-Jeans thermalization/condensation of classical
  light waves {[}11{]} form the lineage in which nonlinear classical
  waves reach statistical equilibrium via recurrences. Our
  \(f^*\approx\) mask phase-space fraction is consistent with the
  equipartition reading of this lineage. However, thermalization under
  an exactly conserved quadratic invariant (the zero closure, drift
  \(10^{-9}\)\ldots{}\(10^{-14}\)) has no counterpart. The Manley-Rowe
  relations {[}12{]} are a distant analogue of ``conservation laws
  constraining conversion paths,'' but those follow from Hamiltonian
  structure, whereas here conservation acts as the axiom that selects
  the vertex.
\item
  \textbf{No counterpart found for rate = C·f²}: in seeded parametric
  amplification the gain rate is set by the pump, and the seed enters
  only as a prefactor. An autocatalytic law in which the growth
  \textbf{rate} itself is proportional to the square of the seed
  fraction could not be found in the surveyed literature (second-order
  autocatalysis in chemical kinetics is a formal similarity only). We
  present it as a novelty candidate and welcome the presentation of
  prior examples (refutation).
\end{enumerate}

\subsection{13. Outlook --- Embedding into the N-Body Model (Verified
Connections)}\label{outlook-embedding-into-the-n-body-model-verified-connections}

The two-body map of this paper has been embedded into the N-body
relational wave system (complete graph, multi-slice register) by a
family of reproducible experiments in the public repository. We record
only the main points (details in the scripts and commit series {[}9{]}).

\begin{enumerate}
\def\labelenumi{\arabic{enumi}.}
\tightlist
\item
  \textbf{``When does the interaction happen'' is resolved}: the
  interaction acts at all times; the apparent timing of generation is
  produced by two factors (condensation gate \(\times\) seed fraction
  squared). The measured gate ratio of generation rate under
  condensation is 2.74.
\item
  \textbf{The uniquification of this paper is the M=2 special case}: on
  general graphs, closure conservation \(\Leftrightarrow\) \(g_2=-g_1\)
  and pair symmetry of the mediated strengths (the mediated-vertex
  uniqueness theorem).
\item
  \textbf{Transplantation of the ignition law}: the many-body system
  reproduces the power 2.001 with \(C\) constant to 0.4\%.
\item
  \textbf{The micro-mechanism of the fate}: the generation direction is
  a phase condition (a 50/50 lottery under random phases), and net
  matter arises by rectification through the \(C f^2\) autocatalysis
  (80\% of the ensemble grows net) --- the decomposition of this paper's
  ``oscillatory equilibrium.''
\item
  \textbf{Position and motion}: locality lives on the register
  (waveform-position) axis, and products are pinned at their birth
  position. Motion is provided by the harmonic-closure dispersion
  \(\omega_k = k\,\omega_1\) {[}3{]} (shape-invariant translation
  demonstrated in 1D/3D).
\item
  \textbf{Transplantation of the census}: the three criteria ---
  sum-rule exclusivity, hair ledger, simultaneity --- hold in the
  many-body system up to the hair-complete version.
\end{enumerate}

\subsubsection{13.1 Open Problems (working guide --- not claimed in this
paper)}\label{open-problems-working-guide-not-claimed-in-this-paper}

The following are analysis problems opened by the results of this paper;
none of them is claimed here.

\begin{enumerate}
\def\labelenumi{\arabic{enumi}.}
\tightlist
\item
  \textbf{Perturbative derivation of the \(f^2\) law}: derive
  analytically, from the vertex and the initial spectrum, why the
  initial growth rate is the square of the seed fraction. Includes
  identifying the mechanism of the common superlinear gain of Figure
  1(a) (common to all seed fractions --- not explained by the current
  \(f^2\) alone). The closed growth law may involve, besides \(f\), an
  alignment quantity of relative phases (e.g.~the distribution of
  \(\mathrm{Im}(\bar b a)\)).
\item
  \textbf{State-count derivation of the equilibrium value}: can the
  \textasciitilde5\% difference between \(f^*=0.4690\) and the naive bin
  ratio \(0.494\) be explained by the \textbf{measure of the reachable
  phase space} under closure, hair, and momentum conservation? If so,
  \(f^*\) is promoted to a statistical-mechanical formula.
\item
  \textbf{Extraction of the generator algebra}: numerically extract the
  generators of the elastic rotation, the inelastic pointwise rotation
  (Section 5.2), and relative-phase operations, and measure whether the
  commutators close. If the elastic and inelastic parts come from one
  generator family, assumption (vi) approaches a derivation. Group
  theory is not input; the operators extracted from the map are
  classified after the fact.
\item
  \textbf{Reduction of the assumptions}: can (ii) pointwise, (iii)
  cubic, (iv) common-phase invariant, (v) no self-scattering be
  re-derived from a single higher principle (local, anonymous,
  relational, lowest-order)?
\end{enumerate}

\subsection{14. What This Paper Does Not Claim;
Falsifiability}\label{what-this-paper-does-not-claim-falsifiability}

This paper does not claim that nature's generative interaction has the
form of Section 4. Assumptions (i)--(vi) can be replaced by other
choices, and all results of this paper are consequences of the
conditional ``if this class is assumed.'' What raises the value of the
assumption is not a proof of correctness but the demonstration of
\textbf{consistency} (no contradiction with the axioms, the existing
elastic operation, and the conserved quantities) and
\textbf{productivity} (the same vertex reproducing the same laws in the
N-body model; Section 13).

Falsification conditions: if any of the predictions of this class ---
sum-rule exclusivity (machine zero of unpredicted bins), the hair ledger
(the machine-derived value of \(q^*\)), the coincidence of the
statistical equilibrium \(f^*\) with the mask phase-space fraction ---
is broken in a more constrained experimental system, the class as a
whole is rejected. If a prior example of the autocatalytic law of
Section 12-6 is presented, the novelty claim is withdrawn.

\subsection{15. Reproducibility}\label{reproducibility}

We follow the series convention (self-contained reproduction within the
series; code of past papers is imported read-only after snapshot
verification).

\begin{itemize}
\tightlist
\item
  \textbf{Original}: the Paper 9 execution environment
  \texttt{ab\_invariant\_theta\_toy\_v1} (unmodified; snapshot SHA-256
  90c7b272\ldots, bit-identity confirmed).
\item
  \textbf{Code of this paper} (repository folder
  \texttt{universal\_inelastic\_map\_managed\_v1}, under the series
  folder for dimensional generation structure):

  \begin{itemize}
  \tightlist
  \item
    \texttt{universal\_inelastic\_map\_v1.py} \ldots{} v1 implementation
    + T1 control / T2 conservation / T3 parity / T4 seeded gain / T5 IF
    audit
  \item
    \texttt{universal\_inelastic\_map\_v3.py} \ldots{} final map
    (strength \(=R\)) + dilute reduction and ignition sweep
  \item
    \texttt{run\_ignition\_fate\_rk4\_v2.py} \ldots{} RK4 fate re-run +
    \(C\) sweep (criteria R1--R3 fixed in advance; calibration record)
  \item
    \texttt{run\_ignition\_fate\_exact\_v3.py} \ldots{} verification of
    the closed-form exact solution (Section 5.2) and re-run of fate and
    sweep (criteria E1--E3 fixed in advance; all passed)
  \item
    \texttt{run\_pair\_structure\_census\_v3.py} \ldots{} final census
    (with preregistered predictions printed)
  \item
    \texttt{make\_paper\_figures\_v1.py} \ldots{} generation of Figures
    1--3 + census RK4 re-measurement
  \end{itemize}
\item
  \textbf{Main parameters}: \(\chi\) 512 bins \(\times\) \(\eta\) 16,
  pump harmonics \(2\!-\!62\) (even), \(S=8\), RK4 substeps
  \(h_{\max}=0.01\) (historical), closed-form implementation for final
  numbers.
\item
  \textbf{Result JSONs}: \texttt{ignition\_fate\_rk4\_result\_v2.json},
  \texttt{ignition\_fate\_exact\_result\_v3.json},
  \texttt{pair\_structure\_census\_result\_v3.json},
  \texttt{pair\_structure\_census\_rk4\_result\_v1.json}, etc. (all
  numbers in the text agree with these files).
\end{itemize}

\subsection{References}\label{references}

\textbf{Within the series (dependence indicated)}

{[}1{]} N. Kihara, \emph{The Generative Structure of Fermions} (2026),
Zenodo. Concept DOI: 10.5281/zenodo.21766706 (original of the elastic
universal interaction, the mask, and the execution environment; Section
3)

{[}2{]} N. Kihara, \emph{Geometric Rapid Expansion Is Specific to
Unstable Self-Consistent Closures} (2026), Zenodo. Concept DOI:
10.5281/zenodo.21798854 (motivation and the correspondence
``amplification \(\Leftrightarrow\) unstable equilibrium''; Sections 1,
6)

{[}3{]} N. Kihara, \emph{Future-Phase Position-Acceleration Map in an AB
Two-Body Closed Phase System and the Inverse-Square Law via Harmonic
Closure} (2026), Zenodo. DOI: 10.5281/zenodo.21441081 (harmonic-closure
dispersion \(\omega_k=k\,\omega_1\); Section 13-5 only)

{[}4{]} N. Kihara, \emph{Two-Grammar Decomposition of Interaction}
(Anonymous Two-Channel Closed Wave System Trilogy, Part I) (2026),
Zenodo. Concept DOI: 10.5281/zenodo.21763995 (identification of the
transmission rate with the elementary charge; motivation of assumption
(vi) only)

\textbf{External}

{[}5{]} C. R. Menyuk, ``Nonlinear pulse propagation in birefringent
optical fibers,'' IEEE J. Quantum Electron. \textbf{23}, 174 (1987)

{[}6{]} G. P. Agrawal, \emph{Nonlinear Fiber Optics}, 6th ed.~(Academic
Press, 2019)

{[}7{]} P. W. Milonni, ``Semiclassical and quantum-electrodynamical
approaches in nonrelativistic radiation theory,'' Phys.
Rep.~\textbf{25}, 1 (1976)

{[}8{]} P. B. Greene, L. Kofman, ``Preheating of fermions,'' Phys. Lett.
B \textbf{448}, 6 (1999); ``Theory of fermionic preheating,'' Phys.
Rev.~D \textbf{62}, 123516 (2000); L. Kofman, A. Linde, A. A.
Starobinsky, ``Towards the theory of reheating after inflation,'' Phys.
Rev.~D \textbf{56}, 3258 (1997)

{[}9{]} This repository, folder
\texttt{universal\_interaction\_manybody\_connection\_v1} (design
policy, the mediated-vertex uniqueness theorem, stages 1--3, GENESIS
v1--v3, the 3D demo --- all scripts, result JSONs, and the commit
series)

{[}10{]} E. Fermi, J. Pasta, S. Ulam (with M. Tsingou), ``Studies of
nonlinear problems I,'' Los Alamos Report LA-1940 (1955)

{[}11{]} K. Baudin \emph{et al.}, ``Classical Rayleigh-Jeans
condensation of light waves,'' Phys. Rev.~Lett. \textbf{125}, 244101
(2020)

{[}12{]} J. M. Manley, H. E. Rowe, ``Some general properties of
nonlinear elements,'' Proc. IRE \textbf{44}, 904 (1956)

\end{document}
